\chapter{The Bias--Variance Decomposition}
\label{ch:biasvar}

This chapter demonstrates how to write mathematics in these notes ---
numbered single-line equations, multi-line derivations, and referring
back to earlier equations and results --- using the classic
bias--variance decomposition as the running example.

\section{Setup}

Let $\theta$ be an unknown parameter and let $\hat\theta = \hat\theta(X_1,\dots,X_n)$
be an estimator of $\theta$ based on a sample $X_1,\dots,X_n$. Since
$\hat\theta$ is a function of random data, it is itself a random variable,
with its own mean and spread.


\begin{definition}[Bias]
\label{def:bias}
The \textbf{bias} of $\hat\theta$ as an estimator of $\theta$ is
\begin{equation}
\label{eq:bias-def}
\operatorname{Bias}(\hat\theta) = \mathbb{E}[\hat\theta] - \theta.
\end{equation}
We say $\hat\theta$ is \textbf{unbiased} if $\operatorname{Bias}(\hat\theta) = 0$,
i.e.\ if $\mathbb{E}[\hat\theta] = \theta$.
\end{definition}

\begin{definition}[Variance]
\label{def:variance}
The \textbf{variance} of $\hat\theta$ is
\begin{equation}
\label{eq:var-def}
\operatorname{Var}(\hat\theta) = \mathbb{E}\big[(\hat\theta - \mathbb{E}[\hat\theta])^2\big].
\end{equation}
\end{definition}

We judge how good $\hat\theta$ is as an estimator of $\theta$ by its
\textbf{mean squared error} (MSE),
\begin{equation}
\label{eq:mse-def}
\operatorname{MSE}(\hat\theta) = \mathbb{E}\big[(\hat\theta - \theta)^2\big].
\end{equation}

\section{The Decomposition}

\begin{proposition}[Bias--variance decomposition]
\label{prop:biasvar}
For any estimator $\hat\theta$ of $\theta$,
\begin{equation}
\label{eq:biasvar}
\operatorname{MSE}(\hat\theta) = \operatorname{Var}(\hat\theta) + \operatorname{Bias}(\hat\theta)^2.
\end{equation}
\end{proposition}

\begin{proof}
Write $\mu = \mathbb{E}[\hat\theta]$ and add and subtract $\mu$ inside the
square in~\eqref{eq:mse-def}:
\begin{align}
\operatorname{MSE}(\hat\theta)
  &= \mathbb{E}\big[(\hat\theta - \theta)^2\big] \notag\\
  &= \mathbb{E}\big[\big((\hat\theta - \mu) + (\mu - \theta)\big)^2\big] \notag\\
  &= \mathbb{E}\big[(\hat\theta - \mu)^2\big]
     + 2(\mu - \theta)\,\mathbb{E}[\hat\theta - \mu]
     + (\mu - \theta)^2. \label{eq:biasvar-expand}
\end{align}
The middle term in \eqref{eq:biasvar-expand} vanishes, because
$\mathbb{E}[\hat\theta - \mu] = \mathbb{E}[\hat\theta] - \mu = 0$ by
definition of $\mu$. The first term is $\operatorname{Var}(\hat\theta)$
by \eqref{eq:var-def}, and the last term is $\operatorname{Bias}(\hat\theta)^2$
by \eqref{eq:bias-def}, since $\mu - \theta = \operatorname{Bias}(\hat\theta)$.
This gives \eqref{eq:biasvar}.
\end{proof}

\begin{remark}
\label{rem:tradeoff}
Proposition~\ref{prop:biasvar} is the reason for the name
\emph{bias--variance tradeoff}: an estimator with zero bias can still
have large MSE if its variance is large, and an estimator with small
variance can still have large MSE if it is badly biased. Choosing a good
estimator (or, later in the course, a good model) usually means trading
some bias for a larger reduction in variance, or vice versa, rather than
minimizing either term alone.
\end{remark}

\section{A Worked Example: Estimating a Variance}

\begin{example}
\label{ex:sample-var}
Let $X_1,\dots,X_n$ be i.i.d.\ with mean $\mu$ and variance $\sigma^2$,
and let $\bar X = \tfrac{1}{n}\sum_{i=1}^n X_i$. Consider the two
estimators of $\sigma^2$,
\begin{equation}
\label{eq:s2biased}
\hat\sigma^2_n = \frac{1}{n}\sum_{i=1}^n (X_i - \bar X)^2,
\qquad
S^2 = \frac{1}{n-1}\sum_{i=1}^n (X_i - \bar X)^2.
\end{equation}
A standard calculation shows
\begin{equation}
\label{eq:s2biased-mean}
\mathbb{E}\big[\hat\sigma^2_n\big] = \frac{n-1}{n}\,\sigma^2,
\end{equation}
so that, by Definition~\ref{def:bias},
\begin{equation}
\operatorname{Bias}(\hat\sigma^2_n) = \frac{n-1}{n}\sigma^2 - \sigma^2 = -\frac{\sigma^2}{n},
\end{equation}
which is nonzero but $\to 0$ as $n \to \infty$: $\hat\sigma^2_n$ is
\emph{biased} but \emph{asymptotically unbiased}. In contrast,
$S^2 = \tfrac{n}{n-1}\hat\sigma^2_n$ satisfies $\mathbb{E}[S^2] = \sigma^2$
exactly, i.e.\ $S^2$ is unbiased for every $n$ -- this is precisely why
the sample variance is conventionally defined with $n-1$ in the
denominator rather than $n$.
\end{example}

\begin{remark}
Example~\ref{ex:sample-var} is a useful sanity check when you meet a new
estimator: an unbiased estimator is not automatically the ``best'' one by
Proposition~\ref{prop:biasvar} -- $S^2$ has larger variance than
$\hat\sigma^2_n$ for normal data, so which one has smaller MSE depends on
$n$ and on how the two terms in \eqref{eq:biasvar} trade off.
\end{remark}
